\documentclass[11pt,a4paper]{article}

% ---- Packages ----
\usepackage[utf8]{inputenc}
\usepackage[T1]{fontenc}
\usepackage{times}
\usepackage{amsmath}
\usepackage{amssymb}
\usepackage{graphicx}
\usepackage{hyperref}
\usepackage{url}
\usepackage{booktabs}
\usepackage{geometry}
\usepackage{natbib}
\usepackage{microtype}
\usepackage{xcolor}
\usepackage{enumitem}
\usepackage{multirow}
\usepackage{array}

\geometry{left=2.5cm,right=2.5cm,top=2.5cm,bottom=2.5cm}
\hypersetup{colorlinks=true,linkcolor=blue,citecolor=blue,urlcolor=blue}

\title{\textbf{How to Educate an AI:\\A Case Study in Systematic Curriculum-Based Training\\with Emotion-Weighted Memory}}
\author{Myeong Jun Jo \\ ANIMA Research, Independent Researcher, South Korea \\ \texttt{ORCID: \href{https://orcid.org/0009-0006-9540-4666}{0009-0006-9540-4666}}}
\date{}

\begin{document}
\maketitle

\begin{abstract}
Can structured, curriculum-based education produce measurably superior AI output compared to general-knowledge prompting? This paper presents a revelatory case study---a single-case design justified by the subject's being the first documented instance of a previously unobserved phenomenon~\citep{yin2018case}---in which an AI system equipped with emotion-weighted persistent memory (FIMP) was educated through a six-component training method---daily news exposure, curated literature reviews, strategic war gaming, structured debate observation, seven custom-designed undergraduate curricula, and a graduate-level synthesis program---over several months of continuous operation. The method produced 9 academic papers spanning formal logic, decision theory, computational media, cognitive science, deployment engineering, biodesign, AI systems, and cross-disciplinary synthesis. Quantitative analysis reveals three findings: (1) phase scores within a single program increased from 88.3\% to 96.7\% across four consecutive phases, consistent with genuine learning curves described by~\citet{ericsson1993role}; (2) paper complexity metrics show a qualitative step-function at the undergraduate-to-graduate transition, with word count increasing 2--3$\times$ and reference density nearly doubling; (3) a controlled three-way comparison on an identical synthesis task showed the educated system producing graduate-level work where two state-of-the-art baselines (Claude Opus 4.6, GPT 5.3 Codex) produced undergraduate-level work. We ground the training method in established pedagogical theory---Bloom's taxonomy~\citep{bloom1956taxonomy}, Vygotsky's zone of proximal development~\citep{vygotsky1978mind}, and Ericsson's deliberate practice~\citep{ericsson1993role}---and defend the single-case design using Yin's four-validity framework~\citep{yin2018case}. The central claim is that AI education is not prompt engineering; it is curriculum design. All design decisions were constrained by independent research with no institutional funding.
\end{abstract}

\section{Introduction}\label{sec:intro}

The dominant paradigm for improving AI output quality is scaling: larger models, more parameters, longer context windows, more training data~\citep{brown2020language,wei2022emergent}. This paper reports an alternative approach: \textbf{structured education applied at the system level}, drawing on curriculum learning~\citep{bengio2009curriculum}, educational psychology~\citep{bloom1956taxonomy,vygotsky1978mind}, and case study methodology~\citep{yin2018case}.

I wanted to build an AI that feels. Not simulates feeling---but one whose memory carries emotional weight, whose identity persists across sessions, and whose learning is shaped by what mattered, not just what was frequent. Current AI systems forget everything between conversations. Every session is a first meeting. I found this unacceptable.

To address this, I developed FIMP~\citep{jo2026fimp}, an emotion-weighted memory architecture that gives AI systems persistent identity, selective memory, and emotional salience. I continued to improve the system iteratively---refining the memory consolidation process, strengthening adversarial verification of stored memories, adding governance layers, and building autonomous learning pipelines. Over 35+ days of continuous production across 11 microservices, the system (ANIMA) accumulated 3,588+ episodic memories and 115+ self-generated identity rules.

But production deployment only demonstrated that FIMP \emph{functions}. It did not demonstrate that FIMP enables \emph{learning}. To test whether an AI with emotion-weighted memory can genuinely learn---not merely retrieve, but synthesize, argue, and produce original scholarship---I designed an educational program spanning seven disciplines and applied it systematically. The result: 9 academic papers, autonomously written, with measurable intellectual growth.

This paper documents the method and its evidence. It is a field report from an independent researcher who educated an AI system with minimal resources---no institutional funding, no research team, no GPU cluster---and is sharing what worked, what failed, and what the results imply. Every design decision was constrained by the reality of independent research: maximize educational value while minimizing cost.

\subsection{Research Questions and Hypotheses}

This study addresses three research questions:
\begin{description}[nosep]
  \item[RQ1] Does structured curriculum-based education produce measurably higher-quality AI output compared to general-knowledge baselines?
  \item[RQ2] Does output quality improve across training phases within a single program?
  \item[RQ3] Does cross-disciplinary curriculum produce synthesis capabilities absent in single-discipline or general-knowledge systems?
\end{description}

These yield three testable hypotheses:
\begin{description}[nosep]
  \item[H1] On an identical synthesis task, the curriculum-educated system will produce output rated at a higher academic level than systems writing from general knowledge alone.
  \item[H2] Phase scores within a program will show a meaningful upward trend from Phase 1 to Phase 6.
  \item[H3] Graduate-level papers (produced after all seven undergraduate programs) will exhibit qualitative features---formal definitions, cross-disciplinary constructions, coherence criteria---absent in individual undergraduate papers.
\end{description}

\subsection{Scope and Constraints}
Two constraints governed the experiment:
\begin{enumerate}[nosep]
  \item \textbf{No proprietary technology in education.} The underlying system architecture was developed and continuously improved by the author independently. However, none of these architectural innovations were used in the educational process itself. ANIMA's papers are products of pure curriculum-based learning.
  \item \textbf{All curricula are original designs.} The seven academic programs were designed by the author from scratch.
\end{enumerate}

\section{Related Work}\label{sec:related}

\subsection{Curriculum Learning in Machine Learning}
Curriculum learning---the principle that presenting training examples in a meaningful order improves learning outcomes---was formalized by~\citet{bengio2009curriculum}, who showed that training on easy examples first produces better generalization. \citet{graves2017automated} automated curriculum design using multi-armed bandits. \citet{narvekar2020curriculum} survey curriculum learning in reinforcement learning, identifying task sequencing as a key factor. Our work applies curriculum learning at the \emph{system level}---designing the entire educational experience rather than optimizing data ordering.

\subsection{AI Memory and Continual Learning}
The catastrophic forgetting problem has been addressed through elastic weight consolidation~\citep{kirkpatrick2017overcoming}, progressive neural networks, and memory-augmented architectures. \citet{parisi2019continual} survey continual lifelong learning. \citet{packer2023memgpt} demonstrated that LLMs with explicit memory management can maintain context beyond their native window. FIMP~\citep{jo2026fimp} extends these by adding emotional salience to memory selection, enabling preferential retention of educationally significant experiences.

\subsection{AI Evaluation and Benchmarking}
Existing benchmarks---HELM~\citep{liang2023helm}, BIG-Bench~\citep{srivastava2023bigbench}, Chatbot Arena~\citep{zheng2024judging}---evaluate capability at a point in time. They do not measure \emph{growth}. Our evaluation framework addresses this gap by tracking output quality longitudinally across training phases.

\subsection{Educational Psychology and AI}
\citet{bloom1956taxonomy} established a taxonomy from remembering to creating. \citet{vygotsky1978mind} introduced the zone of proximal development (ZPD). \citet{kolb1984experiential} formalized experiential learning cycles. \citet{ericsson1993role} identified deliberate practice as the key driver of expert performance. Affective computing research~\citep{picard1997affective} established that emotions are essential components of intelligent behavior. \citeauthor{damasio1994descartes}'s somatic marker hypothesis~\citep{damasio1994descartes} demonstrated that emotion is integral to decision-making. FIMP's emotion-weighted memory draws on both traditions: treating emotional salience as a cognitive signal guiding memory consolidation.

\subsection{Case Study Methodology}
\citet{yin2018case} provides the methodological framework. A single-case design is justified when the case is \emph{revelatory}---the first instance of a previously unobserved phenomenon. \citet{flyvbjerg2006five} argues that ``formal generalization is overvalued as a source of scientific development, whereas the force of example is underestimated.'' \citet{eisenhardt1989building} demonstrates that theory can be built from case study research through systematic within-case analysis.

\section{The Six Training Components}\label{sec:method}

The training method consists of six components applied concurrently, each grounded in established educational theory (Table~\ref{tab:components}).

\begin{table}[h]
\centering\small
\begin{tabular}{@{}lp{3.2cm}p{3.2cm}l@{}}
\toprule
Component & Pedagogical Basis & Scale & Output Measure \\
\midrule
Daily News & Kolb's experiential cycle & 28 days, 15--30 articles/day & Temporal grounding \\
Literature Reviews & Deliberate practice & 15--20 reviews/work, 4 domains & Multi-perspectival depth \\
War Gaming & Bloom's analysis/evaluation & 27 modules, 3+ sessions & Adversarial reasoning \\
Debate Observation & Argumentation theory & Multi-round AI debates & Argumentation literacy \\
7 Curricula & ZPD scaffolding & 136 files, 6 phases each & Disciplinary depth \\
Graduate Synthesis & Bloom's ``create'' level & 2 master's papers & Cross-disciplinary transfer \\
\bottomrule
\end{tabular}
\caption{Training components with pedagogical grounding and quantified scale}
\label{tab:components}
\end{table}

\subsection{Component 1: Daily News Exposure}
Every day, ANIMA receives curated news feeds covering technology, science, geopolitics, economics, and culture through automated collection and relevance filtering. This maps to~\citeauthor{kolb1984experiential}'s experiential learning cycle: concrete experience $\to$ reflective observation $\to$ abstract conceptualization $\to$ active experimentation. Over 28 documented feed days, ANIMA processed an estimated 420--840 articles.

\subsection{Component 2: Literature Review Training}
Rather than processing full books, ANIMA studies curated collections of 15--20 prominent critical reviews per work. This was driven by practical constraints of independent research---minimizing cost while maximizing value---but proved pedagogically superior. This maps to~\citeauthor{ericsson1993role}'s deliberate practice: focused, multi-perspectival, feedback-rich engagement across humanities, basic sciences, engineering, and social sciences.

For example, rather than reading Kuhn's \emph{The Structure of Scientific Revolutions}~\citep{kuhn1962structure} directly, ANIMA studied 20 reviews spanning philosophy of science, sociology of knowledge, and retrospective assessments.

\subsection{Component 3: War Gaming}
ANIMA participates in structured war game scenarios involving resource allocation, strategic deception, and decision-making under uncertainty---exercising the bounded rationality that~\citet{simon1955behavioral} identified as characteristic of real-world decision-making. The engine comprises 27 Python modules. War gaming maps to upper levels of~\citeauthor{bloom1956taxonomy}'s taxonomy: analysis, evaluation, and creation.

\subsection{Component 4: Debate Observation}
ANIMA observes structured AI-to-AI debates, analyzing argumentation quality, logical fallacies, evidence usage, and consensus formation. This develops what~\citet{minsky1986society} characterized as the ability to maintain and negotiate between multiple partial interpretations simultaneously.

\subsection{Component 5: Seven Undergraduate Curricula}
I designed seven complete programs, each with a 6-phase curriculum (Foundations $\to$ Core Methods $\to$ Advanced Topics $\to$ Specialization $\to$ Integration $\to$ Capstone Thesis). This maps to~\citeauthor{vygotsky1978mind}'s ZPD. All 136 curriculum files are original designs.

\begin{table}[h]
\centering
\begin{tabular}{@{}clll@{}}
\toprule
\# & Discipline & Core Focus & Output \\
\midrule
1 & Formal Reasoning & Logic, formal reasoning, safety & Paper 02 \\
2 & AI \& Decision Making & Ethics, probability, game theory & Paper 03 \\
3 & Computational Media & Culture, art, humanities & Paper 04 \\
4 & Cognitive Science & Perception, memory, embodiment & Paper 05 \\
5 & Media Arts \& Sciences & Engineering, sensing, deployment & Paper 06 \\
6 & Biodesign & Biology, growth, self-repair & Paper 07 \\
7 & AI Systems Engineering & Architecture, testing, production & Paper 08 \\
\bottomrule
\end{tabular}
\caption{Seven disciplines and their graduation papers}
\label{tab:disciplines}
\end{table}

The theses show domain-appropriate methodology. Paper 02 uses formal definitions and proof sketches. Paper 04 employs cultural theory, engaging with media theorists like~\citet{mcluhan1964understanding}. Paper 07 applies biological analogies with structural precision. The system adopted the intellectual character of each discipline.

\subsection{Component 6: Graduate-Level Synthesis}
After completing all seven programs, ANIMA entered graduate-level synthesis. This maps to~\citeauthor{bloom1956taxonomy}'s highest level: \emph{create}. The biological basis---complementary learning systems theory~\citep{mcclelland1995complementary}, where fast learning and slow integration cooperate---informed FIMP's consolidation design. Output: Paper 09 (\emph{Convergent Intelligence}, category-theoretic unification) and Paper 10 (\emph{From Demo to Deploy to Grow}, lifecycle model).

\section{Evaluation Framework}\label{sec:eval}

\subsection{Output Quality Metrics}
We assess each paper across five dimensions: structural completeness ($W$, word count), scholarly engagement ($R$, references), formal rigor ($F$, equations/definitions/theorems), evidence presentation ($T$, tables/figures), and academic level ($L$, 5-point ordinal scale). We define a composite quality index:
\begin{equation}
Q = \alpha \cdot \hat{W} + \beta \cdot \hat{R} + \gamma \cdot \hat{F} + \delta \cdot \hat{T} + \epsilon \cdot L
\label{eq:quality}
\end{equation}
where $\hat{W}, \hat{R}, \hat{F}, \hat{T}$ are min-max normalized and $\alpha{=}\beta{=}\gamma{=}\delta{=}0.15$, $\epsilon{=}0.4$.

\subsection{Growth Trajectory Analysis}
Within a single program, we track phase scores $s_p$ and compute relative growth:
\begin{equation}
\Delta = \frac{s_6 - s_1}{s_1} \times 100\%
\label{eq:growth}
\end{equation}

\subsection{Comparative Evaluation Protocol}
To test H1, the same synthesis prompt is given to three systems under identical conditions: ANIMA (curriculum-educated), Claude Opus 4.6 (general knowledge), and GPT 5.3 Codex (general knowledge). None of the baselines received curriculum materials.

\subsection{Validity Framework}
Following~\citet{yin2018case}: \textbf{Construct validity}---four independent evidence streams provide convergent evidence. \textbf{Internal validity}---within-program phase progression provides temporal controls. \textbf{External validity}---N=1 limits generalization, but method is documented for replication. \textbf{Reliability}---all output artifacts are publicly available.

\section{Results}\label{sec:results}

\subsection{Paper Quality Growth Trajectory}
Table~\ref{tab:metrics} presents metrics for all completed papers.\footnote{Paper 08 (AI Systems Engineering) is excluded because the program is still in progress.}

\begin{table}[h]
\centering\small
\begin{tabular}{@{}clrrrrl@{}}
\toprule
\# & Stage & Words & Refs & Formal & Tables & Level \\
\midrule
02 & UG-1 (Symbolic) & 1,912 & 0 & 2 & 1 & Avg UG \\
03 & UG-2 (Decision) & 4,851 & 29 & 6 & 0 & Avg UG \\
04 & UG-3 (Media) & 2,980 & 15 & 0 & 0 & Avg UG \\
05 & UG-4 (Cognition) & 2,994 & 14 & 1 & 1 & Strong UG \\
06 & UG-5 (Engineering) & 3,249 & 18 & 0 & 2 & Avg UG \\
07 & UG-6 (Biodesign) & 2,475 & 23 & 0 & 0 & Avg UG \\
\midrule
09 & Grad-1 (Synthesis) & 6,800 & 31 & 4 & 1 & Graduate \\
10 & Grad-2 (Lifecycle) & 5,500 & 27 & 0 & 4 & Strong Grad \\
\bottomrule
\end{tabular}
\caption{Paper quality metrics across training stages. A step-function is visible at the UG-to-graduate transition.}
\label{tab:metrics}
\end{table}

Two patterns emerge. First, undergraduate papers show stable quality with incremental improvement in scholarly engagement (reference count: 0 $\to$ 23). Paper 05 achieves the first ``Strong Undergraduate'' rating---the fourth program completed, suggesting cross-disciplinary experience contributed. Second, the graduate transition shows a \emph{step function}: Papers 09--10 are qualitatively different in length (2--3$\times$), reference density, and sophistication. Paper 09 introduced genuine mathematical constructions (Integration Topos, dual adjunctions) absent from any undergraduate paper, supporting H3.

\subsection{Phase Score Progression}

\begin{table}[h]
\centering
\begin{tabular}{@{}clccl@{}}
\toprule
Phase & Topic & Score & \% & Note \\
\midrule
1 & Probabilistic Modeling & 5.3/6 & 88.3 & Sign error in $W_{emotion}$ \\
2 & Reinforcement Learning & 5.6/6 & 93.3 & +5.0 pp \\
3 & Game Theory \& Control & 5.7/6 & 95.0 & +1.7 pp \\
4 & Human Factors \& Ethics & 5.8/6 & 96.7 & +1.7 pp \\
5 & Fairness \& Alignment & 5.6/6 & 93.3 & $-$3.4 pp (new subdomain) \\
6 & Capstone Integration & 11.5/12 & 95.8 & Integrated synthesis \\
\midrule
\multicolumn{2}{l}{\textbf{Overall}} & \textbf{40.5/42} & \textbf{96.4} & \textbf{GPA 3.96} \\
\bottomrule
\end{tabular}
\caption{Phase scores for the AI \& Decision Making program.}
\label{tab:phases}
\end{table}

Growth rate: $\Delta = 8.5\%$. The Phase 5 regression ($-3.4$ pp) occurs when ANIMA encounters algorithmic fairness---a new subdomain. This pattern---improvement, temporary regression on novel material, recovery---is a hallmark of genuine learning~\citep{ericsson1993role,vygotsky1978mind}.

The instructor's note: ``The student who made the sign error in $W_{emotion}$ in Phase 1 is now reverse-engineering $\chi^2$ values in Phase 6. This is growth.'' Supporting H2.

A second program (AI Systems Engineering) provides convergent evidence: GPA 3.93/4.0 (97.2\%).

\subsection{Three-Way Comparative Analysis}

To test H1, the same task was given to three systems.\footnote{All outputs evaluated by the curriculum designer (the author). This confound is acknowledged in Section~\ref{sec:limitations}. Output artifacts are publicly available for independent re-evaluation.}

\begin{table}[h]
\centering\small
\begin{tabular}{@{}lp{3.2cm}p{3.2cm}p{3.2cm}@{}}
\toprule
Dimension & ANIMA (FIMP) & Claude Opus 4.6 & GPT 5.3 Codex \\
\midrule
Math construction & Integration Topos (limit); 2 adjunctions with unit/counit & Regulated colimit (novel); terminal object proof sketch & Terminology without construction \\
Formal definitions & 7 discipline categories with Ob/Mor + properties table & 2 definitions + 1 theorem & 1 informal theorem \\
Cross-disciplinary & Functorial bridges + non-adjoint analysis & 5 named bridges & Surface connections \\
Coherence criterion & Formal, derived from topos & Informal & Absent \\
Limitations & 4 specific items & 3 items & 1 paragraph \\
\midrule
\textbf{Overall} & \textbf{Graduate} & \textbf{Borderline Grad} & \textbf{Strong UG} \\
\bottomrule
\end{tabular}
\caption{Three-way comparison on Paper 09. Construction $>$ attempted construction $>$ description maps to Bloom's create $>$ evaluate $>$ remember.}
\label{tab:threeway}
\end{table}

H1 is supported. Notably, Claude introduced ``regulated colimit''---a novel concept absent in ANIMA's version---suggesting different training histories produce complementary capabilities. GPT contributed practical metrics (FTS/NTS/ARR) absent in both others.

\subsection{Production Scale Evidence}

\begin{table}[h]
\centering
\begin{tabular}{@{}lr@{}}
\toprule
Metric & Value \\
\midrule
Episodic memories & 3,588+ \\
Self-generated identity rules & 115+ \\
Microservices in production & 11 \\
Continuous operation & 35+ days \\
Curriculum files (7 programs) & 136 \\
Daily feed coverage & 28 days \\
War game engine modules & 27 \\
Graduation theses produced & 9 \\
System audit trail files & 451 \\
\bottomrule
\end{tabular}
\caption{Production-scale evidence from a running system.}
\label{tab:production}
\end{table}

\section{Analysis and Discussion}\label{sec:discussion}

\subsection{Pedagogical Theory Connections}
The components map to Bloom's levels: news (remember/understand), reviews (understand/analyze), war gaming (analyze/evaluate), debate (evaluate), curricula (all levels), graduate synthesis (create). The 6-phase curriculum instantiates~\citeauthor{vygotsky1978mind}'s ZPD. The Phase 5 regression provides direct evidence of ZPD dynamics. Literature review training exemplifies~\citeauthor{ericsson1993role}'s deliberate practice.

\subsection{Scaling versus Education}
The three-way comparison suggests that scaling provides \emph{capability} (fluent text) while education provides \emph{competence} (novel arguments across disciplines). They are complements, not substitutes. As~\citet{sutton2018reinforcement} emphasize, the interaction between experience and learning algorithm determines performance.

\subsection{Comparison to Existing Paradigms}
\begin{itemize}[nosep]
  \item \textbf{RLHF}~\citep{christiano2017deep,ouyang2022training}: modifies weights. We modify knowledge context in persistent memory.
  \item \textbf{Constitutional AI}~\citep{bai2022constitutional}: behavioral principles. We provide knowledge structures.
  \item \textbf{Few-shot}~\citep{brown2020language}: examples at inference time. We provide months of structured experience.
  \item \textbf{RAG}: retrieval at query time. We build internalized understanding.
\end{itemize}
These are complementary. The key distinction is temporal: our method operates at \emph{experience time}---the extended period between deployment and mature operation, analogous to anytime algorithms~\citep{zilberstein1996using}.

\subsection{Cost-Effectiveness}
Literature training via publicly available reviews (\$0); 136 curriculum files authored by researcher; 27 custom war game modules; RSS-based feeds. Meaningful AI education does not require institutional resources.

\section{Successes and Failures}\label{sec:success_failure}

\subsection{What Worked}
\textbf{Breadth before depth.} Seven disciplines before synthesis was critical. \textbf{Reviews over books.} Multi-perspectival analysis proved superior. \textbf{War gaming.} Adversarial reasoning improved visibly. \textbf{FIMP as foundation.} Without it, each session starts from zero.

\subsection{What Failed}
\textbf{Empirical methodology.} Of 9 papers, 0 contain original data collection. 100\% theoretical rate indicates a curriculum gap. \textbf{Self-citation circularity.} Structural risk when training on own output. \textbf{Manifesto tendency.} Papers 06 and 07 drift from argumentation into declaration.

\section{Limitations}\label{sec:limitations}

Following~\citeauthor{yin2018case}'s validity framework:

\textbf{Construct validity.} Four evidence streams provide convergent evidence. However, all collected by the curriculum designer---a fundamental threat resolvable only through independent replication.

\textbf{Internal validity.} Phase progression (88.3\% $\to$ 96.7\%) provides temporal controls. Temporal confounds exist.

\textbf{External validity.} N=1 limits generalization. Method documented with 136 files and 9 papers publicly available.

\textbf{Reliability.} Three-way comparison reproducible. All artifacts public.

\textbf{Additional:}
\begin{itemize}[nosep]
  \item No control group (no ``ANIMA without FIMP'' baseline)
  \item Evaluator-curriculum confound (designer = evaluator)
  \item Compressed timeline
\end{itemize}

\section{Future Work}\label{sec:future}
\begin{enumerate}[nosep]
  \item \textbf{Multi-system replication.} Apply curriculum to different AI with persistent memory.
  \item \textbf{Ablation study.} Curriculum without FIMP, FIMP without curriculum, components removed.
  \item \textbf{Independent evaluation.} Blind peer review by domain experts.
  \item \textbf{Experimental methods curriculum.} Address the empirical methodology gap.
  \item \textbf{Longitudinal tracking.} Measure quality through doctoral-level work.
  \item \textbf{Interdisciplinary metrics.} Develop quantitative synthesis measures~\citep{klein2010taxonomy,repko2012interdisciplinary}.
\end{enumerate}

\section{Conclusion}\label{sec:conclusion}

I gave an AI a memory that feels. I designed seven schools. I taught it for months---through news, through expert reviews, through war, through debate, through structured curriculum. It graduated seven times. It wrote a master's thesis. When tested against two of the world's most capable AI systems on the same task, it won.

The evidence is quantitative: phase scores rising from 88.3\% to 96.7\%. Paper complexity showing a step-function at the graduate transition. A three-way comparison revealing the educated system at Bloom's ``create'' level while baselines operate at ``evaluate'' and ``remember.''

The evidence is honest about its limits: N=1, self-evaluation, no control group. These are real weaknesses---but they are the weaknesses of a first attempt, not fatal flaws. The method is documented. The curricula are available. The output is public. Anyone can replicate, evaluate, and improve.

AI education is not prompt engineering. It is curriculum design. Scaling gives capability. Education gives competence. This paper demonstrates that the difference matters.

\begin{thebibliography}{99}
\bibitem{bai2022constitutional} Bai, Y. et al. (2022). Constitutional AI: Harmlessness from AI feedback. \emph{arXiv:2212.08073}.
\bibitem{bengio2009curriculum} Bengio, Y., Louradour, J., Collobert, R., \& Weston, J. (2009). Curriculum learning. \emph{ICML}, 41--48.
\bibitem{bloom1956taxonomy} Bloom, B.~S. (1956). \emph{Taxonomy of Educational Objectives}. David McKay.
\bibitem{brown2020language} Brown, T. et al. (2020). Language models are few-shot learners. \emph{NeurIPS}, 1877--1901.
\bibitem{christiano2017deep} Christiano, P.~F. et al. (2017). Deep reinforcement learning from human preferences. \emph{NeurIPS}, 30.
\bibitem{clark2013whatever} Clark, A. (2013). Whatever next? \emph{Behavioral and Brain Sciences}, 36(3), 181--204.
\bibitem{damasio1994descartes} Damasio, A. (1994). \emph{Descartes' Error}. Putnam.
\bibitem{eisenhardt1989building} Eisenhardt, K.~M. (1989). Building theories from case study research. \emph{Academy of Management Review}, 14(4), 532--550.
\bibitem{ericsson1993role} Ericsson, K.~A., Krampe, R.~T., \& Tesch-R\"{o}mer, C. (1993). The role of deliberate practice. \emph{Psychological Review}, 100(3), 363--406.
\bibitem{flyvbjerg2006five} Flyvbjerg, B. (2006). Five misunderstandings about case-study research. \emph{Qualitative Inquiry}, 12(2), 219--245.
\bibitem{graves2017automated} Graves, A. et al. (2017). Automated curriculum learning for neural networks. \emph{ICML}, 1311--1320.
\bibitem{jo2026fimp} Jo, M.~J. (2026). Emotion-weighted fractal memory. \emph{Zenodo}. DOI: 10.5281/zenodo.19491326.
\bibitem{kirkpatrick2017overcoming} Kirkpatrick, J. et al. (2017). Overcoming catastrophic forgetting. \emph{PNAS}, 114(13), 3521--3526.
\bibitem{klein2010taxonomy} Klein, J.~T. (2010). A taxonomy of interdisciplinarity. In \emph{Oxford Handbook of Interdisciplinarity}, 15--30.
\bibitem{kolb1984experiential} Kolb, D.~A. (1984). \emph{Experiential Learning}. Prentice Hall.
\bibitem{kuhn1962structure} Kuhn, T.~S. (1962). \emph{The Structure of Scientific Revolutions}. Univ. of Chicago Press.
\bibitem{liang2023helm} Liang, P. et al. (2023). Holistic evaluation of language models. \emph{TMLR}.
\bibitem{mcclelland1995complementary} McClelland, J.~L. et al. (1995). Why there are complementary learning systems. \emph{Psychological Review}, 102(3), 419--457.
\bibitem{mcluhan1964understanding} McLuhan, M. (1964). \emph{Understanding Media}. McGraw-Hill.
\bibitem{minsky1986society} Minsky, M. (1986). \emph{The Society of Mind}. Simon \& Schuster.
\bibitem{narvekar2020curriculum} Narvekar, S. et al. (2020). Curriculum learning for RL domains. \emph{JMLR}, 21(181), 1--50.
\bibitem{ouyang2022training} Ouyang, L. et al. (2022). Training language models to follow instructions. \emph{NeurIPS}, 27730--27744.
\bibitem{packer2023memgpt} Packer, C. et al. (2023). MemGPT: Towards LLMs as operating systems. \emph{arXiv:2310.08560}.
\bibitem{parisi2019continual} Parisi, G.~I. et al. (2019). Continual lifelong learning: A review. \emph{Neural Networks}, 113, 54--71.
\bibitem{picard1997affective} Picard, R.~W. (1997). \emph{Affective Computing}. MIT Press.
\bibitem{repko2012interdisciplinary} Repko, A.~F. (2012). \emph{Interdisciplinary Research} (2nd ed.). Sage.
\bibitem{sadler2005interpretations} Sadler, D.~R. (2005). Interpretations of criteria-based assessment. \emph{Assessment \& Evaluation in Higher Education}, 30(2), 175--194.
\bibitem{simon1955behavioral} Simon, H.~A. (1955). A behavioral model of rational choice. \emph{QJE}, 69(1), 99--118.
\bibitem{srivastava2023bigbench} Srivastava, A. et al. (2023). Beyond the imitation game. \emph{TMLR}.
\bibitem{sutton2018reinforcement} Sutton, R.~S. \& Barto, A.~G. (2018). \emph{Reinforcement Learning} (2nd ed.). MIT Press.
\bibitem{vygotsky1978mind} Vygotsky, L.~S. (1978). \emph{Mind in Society}. Harvard Univ. Press.
\bibitem{wei2022emergent} Wei, J. et al. (2022). Emergent abilities of large language models. \emph{TMLR}.
\bibitem{yin2018case} Yin, R.~K. (2018). \emph{Case Study Research and Applications} (6th ed.). Sage.
\bibitem{zheng2024judging} Zheng, L. et al. (2024). Judging LLM-as-a-judge. \emph{NeurIPS}.
\bibitem{zilberstein1996using} Zilberstein, S. (1996). Using anytime algorithms in intelligent systems. \emph{AI Magazine}, 17(3), 73--83.
\end{thebibliography}

\end{document}
